%% 自動生成（tools/mkchecklist.py）。各章の「ふりかえり」の項目を集めたもの
\begingroup\footnotesize
\def\bkcb{\tikz[baseline=-0.5ex]\draw[bkTheme, line width=0.7pt, rounded corners=0.5pt] (0,-0.3em) rectangle (0.8em,0.5em);}
\begin{longtblr}{colspec={Q[c,m,wd=9mm] X[l,m] Q[c,m,wd=10mm] Q[c,m,wd=9mm] Q[c,m,wd=9mm]}, width=\linewidth, rowhead=1, colsep=3pt,
  row{1}={bg=bkVol!14, font=\sffamily\bfseries\scriptsize}, hlines={0.4pt,bkLine}, rowsep=2.3pt}
\R{章}{しょう} & できるようになったこと & ページ & 1\R{回}{かい}め & 2\R{回}{かい}め \\
{\sffamily\bfseries\color{bkTheme}2-1} & 「かぞえる\R{数}{かず}」と「はかる\R{数}{かず}」のちがいを\R{説明}{せつ|めい}できる。 & \pageref{bk:2-1:kazoeru} & \bkcb & \bkcb \\
 & \R{量}{りょう}を「\R{数値}{すう|ち}×\R{単位}{たん|い}」で\R{表}{あらわ}し、\R{単位}{たん|い}をかえて\R{表}{あらわ}せる。 & \pageref{bk:2-1:tani} & \bkcb & \bkcb \\
 & \R{小数}{しょう|すう}の\R{位}{くらい}がわかり、10\R{倍}{ばい}・$\tfrac{1}{10}$ で\R{小数点}{しょう|すう|てん}を\R{動}{うご}かせる。 & \pageref{bk:2-1:kurai} & \bkcb & \bkcb \\
{\sffamily\bfseries\color{bkTheme}2-2} & 0より\R{小}{ちい}さい\R{数}{かず}を\R{負}{ふ}の\R{数}{かず}で\R{表}{あらわ}し、\R{数直線}{すう|ちょく|せん}にかける。 & \pageref{bk:2-2:fu} & \bkcb & \bkcb \\
 & \R{不等号}{ふ|とう|ごう}（$<$・$>$・$\leqq$・$\geqq$）で、\R{数}{かず}の\R{大}{おお}きさをくらべられる。 & \pageref{bk:2-2:futougou} & \bkcb & \bkcb \\
 & \R{正}{せい}の\R{数}{かず}・\R{負}{ふ}の\R{数}{かず}の ＋ － × ÷ の\R{計算}{けい|さん}ができる。 & \pageref{bk:2-2:keisan} & \bkcb & \bkcb \\
{\sffamily\bfseries\color{bkTheme}2-3} & \R{同}{おな}じ\R{種類}{しゅ|るい}の\R{量}{りょう}を、\R{単位}{たん|い}をそろえてたしたりひいたりできる。 & \pageref{bk:2-3:onaji} & \bkcb & \bkcb \\
 & ひき\R{算}{ざん}で\R{変化量}{へん|か|りょう}を\R{求}{もと}め、$\Delta$ を\R{使}{つか}って\R{書}{か}ける。 & \pageref{bk:2-3:henka} & \bkcb & \bkcb \\
 & わり\R{算}{ざん}で\R{速}{はや}さを、かけ\R{算}{ざん}で\R{道}{みち}のりを\R{求}{もと}められる。 & \pageref{bk:2-3:wari} & \bkcb & \bkcb \\
{\sffamily\bfseries\color{bkTheme}2-4} & わり\R{算}{ざん}の\R{答}{こた}えを\R{分数}{ぶん|すう}で\R{書}{か}き、\R{約分}{やく|ぶん}・\R{通分}{つう|ぶん}ができる。 & \pageref{bk:2-4:wari} & \bkcb & \bkcb \\
 & \R{分数}{ぶん|すう}の＋－×÷ができ、わり\R{算}{ざん}で\R{逆数}{ぎゃく|すう}をかけるわけを\R{説明}{せつ|めい}できる。 & \pageref{bk:2-4:naze} & \bkcb & \bkcb \\
 & 0と1のあいだの\R{数}{かず}をかけた・わった\R{答}{こた}えの\R{大}{おお}きさがわかる。 & \pageref{bk:2-4:chiisai} & \bkcb & \bkcb \\
{\sffamily\bfseries\color{bkTheme}2-5} & 2つの\R{量}{りょう}の\R{関係}{かん|けい}を\R{比}{ひ}で\R{表}{あらわ}し、\R{等}{ひと}しい\R{比}{ひ}をつくれる。 & \pageref{bk:2-5:hi} & \bkcb & \bkcb \\
 & \R{比}{ひ}をかけ\R{算}{ざん}でつないで、\R{歯車}{は|ぐるま}の\R{回転}{かい|てん}を\R{求}{もと}められる。 & \pageref{bk:2-5:tsunagu} & \bkcb & \bkcb \\
 & \R{割合}{わり|あい}を\R{百分率}{ひゃく|ぶん|りつ}で\R{表}{あらわ}し、「20\%\R{引}{び}き」の\R{量}{りょう}を\R{求}{もと}められる。 & \pageref{bk:2-5:biki} & \bkcb & \bkcb \\
{\sffamily\bfseries\color{bkTheme}2-6} & \R{最小}{さい|しょう}\R{目}{め}もりの10\R{分}{ぶん}の1まで\R{読}{よ}み、\R{平均}{へい|きん}を\R{求}{もと}められる。 & \pageref{bk:2-6:yomu} & \bkcb & \bkcb \\
 & \R{四捨五入}{し|しゃ|ご|にゅう}で、がい\R{数}{すう}をつくれる。 & \pageref{bk:2-6:gaisuu} & \bkcb & \bkcb \\
 & 「\qty{5.0}{cm}」と「\qty{5}{cm}」のちがいを $\pm$ で\R{説明}{せつ|めい}できる。 & \pageref{bk:2-6:yuukou} & \bkcb & \bkcb \\
{\sffamily\bfseries\color{bkTheme}2-7} & \R{同}{おな}じ\R{数}{かず}を\R{何回}{なん|かい}もかけた\R{数}{かず}を、\R{累乗}{るい|じょう}で\R{書}{か}いて\R{計算}{けい|さん}できる。 & \pageref{bk:2-7:ruijo} & \bkcb & \bkcb \\
 & \R{体積}{たい|せき}の\R{単位}{たん|い}や\R{大}{おお}きな\R{数}{かず}を、10の\R{累乗}{るい|じょう}で\R{表}{あらわ}せる。 & \pageref{bk:2-7:juu} & \bkcb & \bkcb \\
 & \R{指数}{し|すう}のきまり、0\R{乗}{じょう}、マイナスの\R{指数}{し|すう}がわかる。 & \pageref{bk:2-7:kimari} & \bkcb & \bkcb \\
{\sffamily\bfseries\color{bkTheme}2-8} & \R{面積}{めん|せき}から1\R{辺}{ぺん}を、2\R{乗}{じょう}ではさんで\R{求}{もと}められる。 & \pageref{bk:2-8:hasami} & \bkcb & \bkcb \\
 & $\sqrt{\phantom{2}}$ の\R{意味}{い|み}がわかり、\R{平方根}{へい|ほう|こん}を\R{求}{もと}められる。 & \pageref{bk:2-8:root} & \bkcb & \bkcb \\
 & \R{周期}{しゅう|き}が\R{長}{なが}さの $\sqrt{\phantom{2}}$ に\R{比例}{ひ|れい}することを\R{説明}{せつ|めい}できる。 & \pageref{bk:2-8:furiko} & \bkcb & \bkcb \\
{\sffamily\bfseries\color{bkTheme}2-9} & \R{分数}{ぶん|すう}を\R{小数}{しょう|すう}にすると、\R{終}{お}わるかくり\R{返}{かえ}すわけがわかる。 & \pageref{bk:2-9:kurikaesu} & \bkcb & \bkcb \\
 & \R{分数}{ぶん|すう}で\R{書}{か}けない\R{数}{かず}（\R{無理数}{む|り|すう}）があることがわかる。 & \pageref{bk:2-9:muri} & \bkcb & \bkcb \\
 & \R{円周}{えん|しゅう}＝\R{直径}{ちょっ|けい}×$\pi$ で、\R{円周}{えん|しゅう}や\R{直径}{ちょっ|けい}を\R{求}{もと}められる。 & \pageref{bk:2-9:pi} & \bkcb & \bkcb \\
{\sffamily\bfseries\color{bkTheme}2-10} & k・m・{\textmu} を\R{読}{よ}んで、\R{単位}{たん|い}を\R{書}{か}きかえられる。 & \pageref{bk:2-10:settou} & \bkcb & \bkcb \\
 & 「\R{単位}{たん|い}つきの1」で、\R{時間}{じ|かん}や\R{速}{はや}さを\R{換算}{かん|さん}できる。 & \pageref{bk:2-10:kansan} & \bkcb & \bkcb \\
 & \R{単位}{たん|い}を\R{計算}{けい|さん}して、\R{式}{しき}が\R{正}{ただ}しいかを\R{確}{たし}かめられる。 & \pageref{bk:2-10:keisan} & \bkcb & \bkcb \\
{\sffamily\bfseries\color{bkTheme}2-11} & ことばの\R{式}{しき}を、□や\R{文字}{も|じ}を\R{使}{つか}った\R{式}{しき}で\R{表}{あらわ}せる。 & \pageref{bk:2-11:moji} & \bkcb & \bkcb \\
 & \R{計算}{けい|さん}の\R{順番}{じゅん|ばん}のきまりと、そのわけがわかる。 & \pageref{bk:2-11:junban} & \bkcb & \bkcb \\
 & どんな\R{数}{かず}でも\R{成}{な}り\R{立}{た}つきまりを、\R{文字}{も|じ}で1\R{行}{ぎょう}に\R{書}{か}ける。 & \pageref{bk:2-11:kimari} & \bkcb & \bkcb \\
{\sffamily\bfseries\color{bkTheme}2-12} & ×をはぶき、÷を\R{分数}{ぶん|すう}にして、\R{文字式}{も|じ|しき}を\R{短}{みじか}く\R{書}{か}ける。 & \pageref{bk:2-12:yakusoku} & \bkcb & \bkcb \\
 & \R{同類項}{どう|るい|こう}をまとめたり、かっこをはずしたりできる。 & \pageref{bk:2-12:dourui} & \bkcb & \bkcb \\
 & \R{指数}{し|すう}のきまりを、\R{文字}{も|じ}を\R{使}{つか}って\R{書}{か}いて\R{使}{つか}える。 & \pageref{bk:2-12:shisu} & \bkcb & \bkcb \\
{\sffamily\bfseries\color{bkTheme}2-13} & \R{公式}{こう|しき}の\R{文字}{も|じ}に\R{数}{かず}を\R{代入}{だい|にゅう}して、\R{式}{しき}の\R{値}{あたい}を\R{求}{もと}められる。 & \pageref{bk:2-13:dainyu} & \bkcb & \bkcb \\
 & \R{単位}{たん|い}ごと\R{代入}{だい|にゅう}し、\R{単位}{たん|い}をそろえて\R{計算}{けい|さん}できる。 & \pageref{bk:2-13:tani} & \bkcb & \bkcb \\
 & マイナスの\R{数}{かず}・\R{分数}{ぶん|すう}・\R{式}{しき}を、かっこごと\R{代入}{だい|にゅう}できる。 & \pageref{bk:2-13:kakko} & \bkcb & \bkcb \\
{\sffamily\bfseries\color{bkTheme}2-14} & \R{等式}{とう|しき}の\R{性質}{せい|しつ}を、\R{天}{てん}びんで\R{説明}{せつ|めい}できる。 & \pageref{bk:2-14:tenbin} & \bkcb & \bkcb \\
 & $x=vt$ を $t$ や $v$ について\R{解}{と}ける。 & \pageref{bk:2-14:vt} & \bkcb & \bkcb \\
 & \R{移項}{い|こう}を\R{使}{つか}い、1\R{行}{ぎょう}ずつきまりを\R{書}{か}いて\R{変形}{へん|けい}できる。 & \pageref{bk:2-14:ikou} & \bkcb & \bkcb \\
{\sffamily\bfseries\color{bkTheme}2-15} & \R{比例}{ひ|れい}を\R{見}{み}つけ、$y=ax$ の\R{式}{しき}とグラフにできる。 & \pageref{bk:2-15:hirei} & \bkcb & \bkcb \\
 & ばねの $F=kx$、オームの\R{法則}{ほう|そく} $V=RI$ を\R{使}{つか}える。 & \pageref{bk:2-15:bane} & \bkcb & \bkcb \\
 & \R{反比例}{はん|ぴ|れい}を\R{見分}{み|わ}け、2\R{倍}{ばい}にするとどうなるか\R{答}{こた}えられる。 & \pageref{bk:2-15:hanpirei} & \bkcb & \bkcb \\
{\sffamily\bfseries\color{bkTheme}2-16} & \R{方程式}{ほう|てい|しき}をつくり、\R{表}{ひょう}でためして\R{解}{かい}をさがせる。 & \pageref{bk:2-16:houteishiki} & \bkcb & \bkcb \\
 & 1\R{次方程式}{じ|ほう|てい|しき}を1\R{行}{ぎょう}ずつ\R{解}{と}ける。 & \pageref{bk:2-16:tejun} & \bkcb & \bkcb \\
 & \R{表}{ひょう}・グラフ・\R{方程式}{ほう|てい|しき}で\R{解}{と}き、\R{確}{たし}かめられる。 & \pageref{bk:2-16:graph} & \bkcb & \bkcb \\
{\sffamily\bfseries\color{bkTheme}2-17} & 2\R{進数}{しん|すう}と10\R{進数}{しん|すう}を\R{行}{い}き\R{来}{き}できる。 & \pageref{bk:2-17:henkan} & \bkcb & \bkcb \\
 & ビットの\R{数}{かず}から、\R{何}{なん}\R{通}{とお}りの\R{数}{かず}を\R{表}{あらわ}せるかを\R{計算}{けい|さん}できる。 & \pageref{bk:2-17:bit} & \bkcb & \bkcb \\
 & 2\R{進数}{しん|すう}を4けたずつまとめて、16\R{進数}{しん|すう}で\R{書}{か}ける。 & \pageref{bk:2-17:juroku} & \bkcb & \bkcb \\
{\sffamily\bfseries\color{bkTheme}2-18} & 「＝」の\R{式}{しき}を\R{書}{か}き、コピーしてくり\R{返}{かえ}せる。 & \pageref{bk:2-18:shiki} & \bkcb & \bkcb \\
 & ずらしたくない\R{番地}{ばん|ち}に \$ をつけて\R{使}{つか}える。 & \pageref{bk:2-18:dollar} & \bkcb & \bkcb \\
 & \R{上}{うえ}の\R{行}{ぎょう}の\R{結果}{けっ|か}で\R{下}{した}の\R{行}{ぎょう}を\R{計算}{けい|さん}し、グラフにできる。 & \pageref{bk:2-18:kurikaesi} & \bkcb & \bkcb \\
{\sffamily\bfseries\color{bkTheme}2-19} & \R{変数}{へん|すう}と\R{代入}{だい|にゅう}のプログラムの\R{結果}{けっ|か}を\R{予想}{よ|そう}できる。 & \pageref{bk:2-19:dainyu} & \bkcb & \bkcb \\
 & くり\R{返}{かえ}しと「もし～なら」のプログラムを\R{読}{よ}める。 & \pageref{bk:2-19:kurikaesi} & \bkcb & \bkcb \\
 & \R{小数}{しょう|すう}のずれのわけと\R{対策}{たい|さく}を\R{説明}{せつ|めい}できる。 & \pageref{bk:2-19:zure} & \bkcb & \bkcb \\
\end{longtblr}
\endgroup
